\section{Common methodology}
\label{sec:common_methodology}

\subsection{Research question and experimental scope}
The September 30 proposal asks whether selecting a sports-outcome model for
calibration, rather than classification accuracy, improves betting returns.
Aman Dhillon implements the NBA component and Panagiotis Housos the English
Premier League (EPL) component. Both studies connect pre-game statistical
features to outcome probabilities and then to explicit betting rules. A good
classifier, a well-calibrated probability forecast and a profitable betting
strategy are different objects: a profitable bet also requires a sufficiently
favourable price and a viable stake.

The motivating reference is Walsh and Joshi's study of NBA model
selection~\cite{walsh2024}. Its 2025 corrigendum supersedes the original
numerical claims: corrected mean capital returns are $-9.77\%$ for
calibration-based selection and $-26.78\%$ for accuracy-based selection, and
Kelly betting can lose despite calibration~\cite{walsh2025corrigendum}.
Consequently, improved profitability is tested rather than assumed. The NBA
experiment is a replication and rolling extension; the EPL experiment adapts
the pipeline to three outcomes and subsequently adds exploratory team, xG and
player-feature comparisons. Multi-book line shopping and mathematical
arbitrage screening are implemented in the EPL component, not the supplied
NBA notebook's single-book replay.

\subsection{Pre-game information and chronological model selection}
For fixture $i$, let $Y_i$ denote the realised result and $X_i$ the feature
vector constructed from earlier completed games. NBA results are binary,
$Y_i\in\{H,A\}$; EPL results include draws,
$Y_i\in\{H,D,A\}$. Team features describe earlier opponent-relative performance
and prior-season strength. The NBA code applies shifted expanding averages
within each team-season. EPL features include season-to-date and recent form,
with source-coverage controls and, in the extensions, lagged xG and FPL player
aggregates. Target-game scores, appearances and minutes cannot be predictors
of that same game.

The initial experiments use separate accuracy- and calibration-driven
branches, each considering logistic regression, random forests, support vector
machines and multilayer perceptrons. Distribution screening, Spearman
redundancy filtering, sequential forward selection and TPE hyperparameter
search precede the later betting evaluation. These are shared methodological
ideas, not identical settings. All normalisation statistics are estimated in
the appropriate earlier fitting block. Table~\ref{tab:league_method_differences}
records differences that matter for interpreting the results.

\begin{table}[htbp]
\centering
\small
\caption{Shared framework and league-specific implementation.}
\label{tab:league_method_differences}
\begin{tabular}{>{\raggedright\arraybackslash}p{0.20\textwidth}>{\raggedright\arraybackslash}p{0.35\textwidth}>{\raggedright\arraybackslash}p{0.35\textwidth}}
\toprule
Aspect & NBA & Premier League \\
\midrule
Outcome and feature eligibility & Home/away; both teams must have at least
ten earlier season games. & Home/draw/away; missing histories retained with
coverage and fitting-only missing-value handling. \\
First selection chronology & Train 2014/15--2015/16; tune 2016/17; choose
model family on 2017/18; bet 2018/19. & Train 2014/15--2020/21; tune 2021/22;
calibrate 2022/23; select family 2023/24; choose threshold 2024/25; evaluate
2025/26. \\
Calibration treatment & Selection by a 20-bin binary ECE with an occupancy
penalty; no separate held-out temperature fit. & Selection by mean three-class
ECE with ten bins; both branches receive held-out temperature scaling. \\
Extensions & Nine rolling windows reselect features but fit logistic
regression in both branches. & Four rolling evaluation seasons compare fixed
feature groups; the later player-model study adds logistic and boosted trees. \\
Prices and outcome choice & One saved home/away moneyline pair; test home
value first, otherwise away value. & Saved individual-book quotes and a
Bet365 reference; select the allowed outcome with maximum estimated return. \\
Capital allocation & Sequential per-game settlement; fixed 100 or one-eighth
Kelly from 10,000 initial capital. & Same-date opening-cash allocation;
fixed, cash-fraction and capped fractional-Kelly configurations; 100,000
initial capital in the expanded account comparison. \\
\bottomrule
\end{tabular}
\end{table}

Selecting a model for small calibration error is not the same as fitting a
post-processing calibrator. In the first EPL study, a scalar temperature is
fitted on a disjoint season for every candidate in both branches,
$p_{ik}(T)=\exp(z_{ik}/T)/\sum_j\exp(z_{ij}/T)$, where $z_i$ are the saved
model scores~\cite{guo2017calibration}. The NBA notebook instead selects using
its calibration criterion. Neither implementation implies that probability
calibration guarantees value against bookmaker prices.

\subsection{Forecast scores and market benchmarks}
Let $p_{ik}$ be the predicted probability of outcome $k$ and
$y_{ik}=\mathbb{1}\{Y_i=k\}$. Classification accuracy uses the highest-probability
outcome. Probability quality is assessed using log loss and Brier loss:
\begin{equation}
 \operatorname{LL}=-\frac{1}{n}\sum_i\log p_{i,Y_i},\qquad
 \operatorname{Brier}_{\mathrm{EPL}}=\frac{1}{n}\sum_i\sum_{k\in\{H,D,A\}}
 (p_{ik}-y_{ik})^2.
\end{equation}
The NBA binary Brier convention is $n^{-1}\sum_i(p_{iH}-y_{iH})^2$.
Its values therefore should not be ranked directly against the unscaled
three-class EPL Brier score. Log-loss levels also depend on the outcome space
and base rates; model--market differences are interpreted within each league.

For probability bin $b$ of class $k$, calibration error compares observed
frequency $\bar y_{kb}$ with mean forecast $\bar p_{kb}$:
\begin{equation}
 \operatorname{ECE}_k=\sum_b\frac{n_{kb}}{n}
       |\bar y_{kb}-\bar p_{kb}|.
\end{equation}
The EPL implementation averages this quantity over all three classes and
reports ten-bin ECE, with five- and fifteen-bin sensitivity summaries. The NBA
implementation uses the home-win probability and twenty equal-width bins;
it replaces the score by $1$ if fewer than 80\% of bins are occupied. This is
an optimisation penalty, not an observation of literal 100\% miscalibration.
The two ECE scales are not interchangeable. Constant forecasts may have small
unpenalised ECE while offering little match-specific information.

For decimal odds $O_{ik}$, the implied probability is $1/O_{ik}$. A reference
book's overround and proportional margin-normalised benchmark are
\begin{equation}
 A_i=\sum_k O_{ik}^{-1},\qquad
 q_{ik}=\frac{O_{ik}^{-1}}{A_i}.
\end{equation}
The benchmark $q_i$ is a market estimate, not known truth. NBA American
moneylines $m$ are converted to decimal odds as $1+m/100$ for $m>0$, and
$1+100/|m|$ for $m<0$. EPL line shopping uses the largest eligible archived
individual-book quote for each outcome; a separate reference-price replay
recomputes selections and stake sizes. An archived combination with
$\sum_k1/O^*_{ik}<1$ is a mathematical arbitrage candidate only. Without
simultaneous accepted fills, aligned settlement and known costs, it is not an
executed risk-free profit.

\subsection{Bet selection, stake sizing and accounting}
For selected outcome probability $p_i$ and decimal odds $O_i$, estimated
return per unit staked is $e_i=p_iO_i-1$. The NBA rule takes a home bet when
$p_{iH}>1/O_{iH}$, otherwise an away bet when
$1-p_{iH}>1/O_{iA}$; it does not maximise across both qualifying sides.
The principal EPL value rules instead choose the largest admissible $e_i$
and require a declared threshold. Thresholds and filters are reported with
their respective strategies; the first EPL validation-selected rules can
remain entirely in cash.

With stake $s_i$, the realised net payoff and the model-implied positive and
negative cash components are
\begin{align}
 \pi_i&=s_i\{O_i\mathbb{1}(\text{win})-1\},\\
 G_i&=p_i s_i(O_i-1),\qquad L_i=(1-p_i)s_i,
 \qquad G_i-L_i=s_i e_i.
\end{align}
Expected gains and losses use model probabilities at the recorded stakes;
they are distinct from average realised winning gains and losing stakes.
They are reported only where the supplied prediction-and-trade evidence
supports their calculation.

Full Kelly requests fraction $f_i^*=(p_iO_i-1)/(O_i-1)$ of available capital.
Fractional Kelly multiplies this by $\lambda$, with any declared per-bet cap
$c$: $s_i=B\min\{c,\lambda\max(0,f_i^*)\}$.
The NBA notebook uses $\lambda=1/8$ and caps only at available cash. EPL
configurations specify their Kelly fraction, per-bet cap and any total
date-level budget. A fixed stake of 1,000 in a 100,000 account equals 1\%
of \emph{initial} capital, not 1\% of a changing bankroll. EPL same-date
requests are reduced together to fit opening cash before any same-date
winnings settle. The NBA source simulator processes games sequentially, so
its funding convention is retained and disclosed rather than silently
replaced by the EPL convention.

For total profit $P=\sum_i\pi_i$ and turnover $V=\sum_i s_i$, the final report
uses two explicitly different denominators:
\begin{equation}
 \text{stake ROI (yield)}=P/V,\qquad
 \text{capital return}=P/B_0.
\end{equation}
The supplied NBA material calls the second quantity ``ROI''; it is relabelled
as capital return for consistency. Hit ratio is wins divided by settled bets;
trade frequency is bets divided by priced match opportunities. Drawdown is
the largest proportional fall from a prior account peak. Undefined no-trade
ratios remain unavailable. Where a display rebases an account to 100,000,
all cash amounts and stakes are multiplied by the same constant; reported
percentage returns do not change. Separately reset NBA seasons are not
presented as one funded continuous investment account.

\subsection{Interpretation and evidence limits}
The completed historical experiments are comparisons of recorded designs,
not new live trading records. Repeatedly examined EPL extensions are explicitly
exploratory. NBA annual variation and EPL whole-date conditional bootstrap
intervals do not remove feature, model or strategy selection uncertainty.
Bookmaker probabilities are competitive benchmarks, and favourable model
expected value is not independent evidence of a realised edge. Neither a
positive season nor a historical rank establishes a reliable betting policy.
The NBA results below preserve the supplied notebook and report evidence;
the EPL section retains the frozen prediction, trade and cash-flow artifacts.
